\scanpage{27}%
\mspage{12}%
\margin{\handwritten{$\begin{aligned}j&=0 &&\Longleftrightarrow c_4=0\\j&=12^3 &&\Longleftrightarrow c_6=0\end{aligned}$}}%
\[
j=c_4^3/\Delta=1728+c_6^2/\Delta \, . \tag*{(6.2)}
\]

\uline{Proposition} 6.3. \quad \uline{Soient} $E$ \uline{et} $E'$ \uline{deux courbes de genre un sur} $S$, \uline{sans fibre additive} .

(I) $\text{\uline{Isom}}_{\handwritten{S}}(E,E')$ \uline{est représentable, fini et} \mbox{\begingroup\edef\encoutersep{\the\dimexpr.5\pgflinewidth\relax}\setbox0=\hbox{\uline{non ramifié}}\setbox2=\hbox{\inserted{\handwritten{net}}}\edef\encwda{\the\wd0}\edef\enchta{\the\ht0}\edef\encdpa{\the\dp0}\edef\encwdb{\the\wd2}\edef\enchtb{\the\ht2}\edef\encdpb{\the\dp2}\rlap{\smash{\tikz[baseline=0pt]{\begin{scope}[shift={(\the\dimexpr \encwda/2\relax,\the\dimexpr -\encdpa/2+\enchta/2\relax)}]\coordinate (enc-nrone-north-east) at (\the\dimexpr \encoutersep+\encwda/2\relax,\the\dimexpr \encdpa/2+\enchta/2+\encoutersep\relax);\coordinate (enc-nrone-north-west) at (\the\dimexpr -\encoutersep-\encwda/2\relax,\the\dimexpr \encdpa/2+\enchta/2+\encoutersep\relax);\coordinate (enc-nrone-south-east) at (\the\dimexpr \encoutersep+\encwda/2\relax,\the\dimexpr -\encdpa/2-\enchta/2-\encoutersep\relax);\coordinate (enc-nrone-south-west) at (\the\dimexpr -\encoutersep-\encwda/2\relax,\the\dimexpr -\encdpa/2-\enchta/2-\encoutersep\relax);\coordinate (enc-nrone-west) at (\the\dimexpr -\encoutersep-\encwda/2\relax,0pt);\handwritten{\draw[thin,gray] (enc-nrone-north-east) .. controls ([xshift=-1em,yshift=0.8em]enc-nrone-north-east) and ([xshift=-0.6em,yshift=0.7em]enc-nrone-north-west) .. ([xshift=-0.2em]enc-nrone-west) .. controls ([yshift=-0.85em]enc-nrone-south-west) and ([xshift=0.5em,yshift=-0.85em]enc-nrone-south-east) .. (enc-nrone-north-east);}\end{scope}\pgfresetboundingbox\path[use as bounding box] (0pt,\the\dimexpr-\encdpa\relax) rectangle (\encwda,\enchta);}}}\rlap{\kern\the\dimexpr \encoutersep*2+\encwda-1.8em\relax\raisebox{\the\dimexpr \encdpb+\enchta+\encoutersep*2+0.48em\relax}[0pt][0pt]{\copy2}}\copy0\endgroup} \uline{sur} $S$\mbox{\begingroup\edef\encoutersep{\the\dimexpr.5\pgflinewidth\relax}\setbox0=\hbox{.}\setbox2=\hbox{\handwritten{\inserted{\handwritten{,}}}}\edef\encwda{\the\wd0}\edef\enchta{\the\ht0}\edef\encdpa{\the\dp0}\edef\encwdb{\the\wd2}\edef\enchtb{\the\ht2}\edef\encdpb{\the\dp2}\rlap{\smash{\tikz[baseline=0pt]{\begin{scope}[shift={(\the\dimexpr \encwda/2\relax,\the\dimexpr -\encdpa/2+\enchta/2\relax)}]\coordinate (enc-stop-center) at (0pt,0pt);\handwritten{\draw[thin,gray,rotate around={-18:(enc-stop-center)}] (enc-stop-center) ellipse (0.56em and 1.34em);}\end{scope}\pgfresetboundingbox\path[use as bounding box] (0pt,\the\dimexpr-\encdpa\relax) rectangle (\encwda,\enchta);}}}\rlap{\kern\the\dimexpr \encwda/2-\encwdb/2\relax\raisebox{\the\dimexpr -\encdpa-\enchtb-\encoutersep*2-0.1em\relax}[0pt][0pt]{\copy2}}\copy0\endgroup}

(II) \uline{si} $j_E=j_F$, \uline{la projection de} $\text{\uline{Isom}}_{\handwritten{S}}(E,E')$ \uline{sur} $S$ \uline{est surjective}\mbox{\begingroup\edef\encoutersep{\the\dimexpr.5\pgflinewidth\relax}\setbox0=\hbox{.}\setbox2=\hbox{\handwritten{\inserted{\handwritten{,}}}}\edef\encwda{\the\wd0}\edef\enchta{\the\ht0}\edef\encdpa{\the\dp0}\edef\encwdb{\the\wd2}\edef\enchtb{\the\ht2}\edef\encdpb{\the\dp2}\rlap{\smash{\tikz[baseline=0pt]{\begin{scope}[shift={(\the\dimexpr \encwda/2\relax,\the\dimexpr -\encdpa/2+\enchta/2\relax)}]\coordinate (enc-stop-center) at (0pt,0pt);\handwritten{\draw[thin,gray,rotate around={-18:(enc-stop-center)}] (enc-stop-center) ellipse (0.56em and 1.34em);}\end{scope}\pgfresetboundingbox\path[use as bounding box] (0pt,\the\dimexpr-\encdpa\relax) rectangle (\encwda,\enchta);}}}\rlap{\kern\the\dimexpr \encwda/2-\encwdb/2\relax\raisebox{\the\dimexpr -\encdpa-\enchtb-\encoutersep*2-0.1em\relax}[0pt][0pt]{\copy2}}\copy0\endgroup}

(III) \uline{si} $j_E=j_F$ \uline{ne prend pas les valeurs} $0$ \uline{et} $\underset{\inserted{\handwritten{\substack{\Vert\\12^3}}}}{\mbox{\begingroup\edef\encoutersep{\the\dimexpr.5\pgflinewidth\relax}\setbox0=\hbox{$1728$}\edef\encwda{\the\wd0}\edef\enchta{\the\ht0}\edef\encdpa{\the\dp0}\rlap{\smash{\tikz[baseline=0pt]{\begin{scope}[shift={(\the\dimexpr \encwda/2\relax,\the\dimexpr -\encdpa/2+\enchta/2\relax)}]\coordinate (enc-val-south-east) at (\the\dimexpr \encoutersep+\encwda/2\relax,\the\dimexpr -\encdpa/2-\enchta/2-\encoutersep\relax);\coordinate (enc-val-south-west) at (\the\dimexpr -\encoutersep-\encwda/2\relax,\the\dimexpr -\encdpa/2-\enchta/2-\encoutersep\relax);\handwritten{\draw[thin] (enc-val-south-west) .. controls ([yshift=-0.15em]enc-val-south-west) and ([yshift=-0.15em]enc-val-south-east) .. (enc-val-south-east);}\end{scope}\pgfresetboundingbox\path[use as bounding box] (0pt,\the\dimexpr-\encdpa\relax) rectangle (\encwda,\enchta);}}}\copy0\endgroup}}$, $\text{\uline{Isom}}_{\handwritten{S}}(E,E')$ \uline{est un revêtement étale de rang} $2$ \uline{de} $S$ .

Donnons-nous $E$ et $E'$ par des équations (1.3) et $(1.3)'$. Se donner un isomorphisme entre $E$ et $E'$ revient à se donner $u$, $r$, $s$ et $t$, vérifiant les équations (1.6), donc aussi (1.7) et (1.8), $u$ devant de plus être inversible. Par hypothèse, localement sur $S$, $c_4$ ou $\Delta$ est inversible (6.1), et de même $c'_4$ ou $\Delta'$ l'est. Localement sur $S$, les équations (1.8) sont donc des équations de dépendance intégrale pour $u$ et $u^{-1}$, les équations $(1.7)_6$ ou $(1.7)_8$ des équations de dépendance intégrale pour $r$ et enfin $(1.6)_6$ (resp. $(1.6)_2$) est une équation de dépendance intégrale pour $t$ (resp. $s$). Ceci prouve que $\text{\uline{Isom}}(E,E')$ est représentable par un schéma affine fini sur $S$ ; il est même \mbox{\begingroup\edef\encoutersep{\the\dimexpr.5\pgflinewidth\relax}\setbox0=\hbox{non ramifié}\setbox2=\hbox{\inserted{\handwritten{net}}}\edef\encwda{\the\wd0}\edef\enchta{\the\ht0}\edef\encdpa{\the\dp0}\edef\encwdb{\the\wd2}\edef\enchtb{\the\ht2}\edef\encdpb{\the\dp2}\rlap{\smash{\tikz[baseline=0pt]{\begin{scope}[shift={(\the\dimexpr \encwda/2\relax,\the\dimexpr -\encdpa/2+\enchta/2\relax)}]\coordinate (enc-nrtwo-north-east) at (\the\dimexpr \encoutersep+\encwda/2\relax,\the\dimexpr \encdpa/2+\enchta/2+\encoutersep\relax);\coordinate (enc-nrtwo-north-west) at (\the\dimexpr -\encoutersep-\encwda/2\relax,\the\dimexpr \encdpa/2+\enchta/2+\encoutersep\relax);\coordinate (enc-nrtwo-south-east) at (\the\dimexpr \encoutersep+\encwda/2\relax,\the\dimexpr -\encdpa/2-\enchta/2-\encoutersep\relax);\coordinate (enc-nrtwo-south-west) at (\the\dimexpr -\encoutersep-\encwda/2\relax,\the\dimexpr -\encdpa/2-\enchta/2-\encoutersep\relax);\coordinate (enc-nrtwo-west) at (\the\dimexpr -\encoutersep-\encwda/2\relax,0pt);\handwritten{\draw[thin,gray] (enc-nrtwo-north-east) .. controls ([xshift=-1em,yshift=0.8em]enc-nrtwo-north-east) and ([xshift=-0.6em,yshift=0.7em]enc-nrtwo-north-west) .. ([xshift=-0.2em]enc-nrtwo-west) .. controls ([yshift=-0.85em]enc-nrtwo-south-west) and ([xshift=0.5em,yshift=-0.85em]enc-nrtwo-south-east) .. (enc-nrtwo-north-east);}\end{scope}\pgfresetboundingbox\path[use as bounding box] (0pt,\the\dimexpr-\encdpa\relax) rectangle (\encwda,\enchta);}}}\rlap{\kern\the\dimexpr \encoutersep*2+\encwda-1.8em\relax\raisebox{\the\dimexpr \encdpb+\enchta+\encoutersep*2+0.48em\relax}[0pt][0pt]{\copy2}}\copy0\endgroup} car les courbes considérées n'ont pas d'automorphismes infinitésimaux (respectant l'origine).

Lorsque $2$ est inversible sur $S$, l'assertion (III) résulte de l'énoncé plus précis.

\uline{Proposition} 6.4. \quad \uline{Soient} $E$ \uline{et} $E'$ \uline{deux courbes de genre un sur} $S$ ,%
% archive sheet number: 26
